% ============================================================
% Secao 4: Directional Relational Manifold (DRM)
% ============================================================

\section{Directional Relational Manifold (DRM)}
\label{sec:drm}

O \textit{Directional Relational Manifold} (DRM) \'e o substrato
geom\'etrico sobre o qual opera o sistema epist\^emico. Cada token \'e
mapeado para um ponto em um manifold Riemanniano 5D aprendido, onde
dist\^ancias geod\'esicas codificam rela\c{c}\~oes epist\^emicas.

\subsection{Embedding no Manifold}

O mapeamento dos \textit{hidden states} para coordenadas 5D \'e feito por
uma MLP com ativa\c{c}\~ao sigmoid na \'ultima camada:

\begin{equation}
    \bm{c} = \sigma\big(\bm{W}_2 \cdot \text{GELU}(\bm{W}_1 \bh + \bm{b}_1) + \bm{b}_2\big) \in [0, 1]^5
    \label{eq:manifold_emb}
\end{equation}

com $\bm{W}_1 \in \R^{h \times d}$, $\bm{W}_2 \in \R^{5 \times h}$,
onde $h = 5 \times 4 = 20$ (configur\'avel via \texttt{hidden\_factor}),
totalizando {\raise.17ex\hbox{$\scriptstyle\sim$}}16K par\^ametros para $d = 768$.

O embedding tamb\'em computa dist\^ancias L2 a 6 \textbf{pontos \^ancora}
fixos no manifold (verdade, ignor\^ancia, ru\'ido, complexo, obsoleto,
ideal), refinados via uma pequena camada de proje\c{c}\~ao.

As 5 dimens\~oes capturam aspectos epist\^emicos distintos:

\begin{table}[H]
\centering
\caption{Interpreta\c{c}\~ao das 5 dimens\~oes do manifold.}
\label{tab:dims_5d}
\begin{tabular}{cl}
\toprule
\textbf{Dim.} & \textbf{Interpreta\c{c}\~ao} \\
\midrule
$c_0$ & Certeza factual \\
$c_1$ & Consist\^encia l\'ogica \\
$c_2$ & Contexto / \textit{grounding} \\
$c_3$ & Calibra\c{c}\~ao de confian\c{c}a \\
$c_4$ & Profundidade epist\^emica \\
\bottomrule
\end{tabular}
\end{table}

\subsection{Tensor M\'etrico}

A geometria do manifold \'e definida por um tensor m\'etrico
$\bG \in \R^{5 \times 5}$ que deve ser \textbf{sim\'etrico positivo-definido}
(SPD). Constru\'imos $\bG$ a partir de uma matriz triangular inferior
aprendida $\bm{L}$:

\begin{equation}
    \bG = \bm{L}\bm{L}^\top + \epsilon \bm{I}_5
    \label{eq:metric_spd}
\end{equation}

onde $\bm{L}$ \'e parametrizado por $5 \times (5+1) / 2 = 15$ par\^ametros
livres e $\epsilon = 10^{-6}$ garante estabilidade num\'erica. A
decomposi\c{c}\~ao de Cholesky garante que $\bG$ \'e SPD por constru\c{c}\~ao.

\begin{proposition}[Positividade do Tensor M\'etrico]
Para qualquer $\bm{L} \in \R^{5 \times 5}$ triangular inferior com diagonal
n\~ao-nula, $\bG = \bm{L}\bm{L}^\top + \epsilon \bm{I}$ \'e SPD.
\end{proposition}

\begin{proof}
Para qualquer $\bm{v} \neq 0$:
$\bm{v}^\top \bG \bm{v} = \bm{v}^\top \bm{L}\bm{L}^\top \bm{v} + \epsilon \|\bm{v}\|^2
= \|\bm{L}^\top \bm{v}\|^2 + \epsilon \|\bm{v}\|^2 > 0$.
\end{proof}

\subsection{Dist\^ancia Geod\'esica}

A dist\^ancia geod\'esica entre dois pontos $\bm{c}_i, \bm{c}_j$ no manifold \'e
aproximada pela dist\^ancia de Mahalanobis induzida por $\bG$:

\begin{equation}
    d_{\bG}(\bm{c}_i, \bm{c}_j) = \sqrt{(\bm{c}_i - \bm{c}_j)^\top \bG\, (\bm{c}_i - \bm{c}_j)}
    \label{eq:geodesic}
\end{equation}

Esta \'e uma aproxima\c{c}\~ao de primeira ordem da dist\^ancia geod\'esica
verdadeira, v\'alida quando a curvatura \'e localmente uniforme. Para o caso
geral, a geod\'esica exigiria integra\c{c}\~ao da equa\c{c}\~ao diferencial:

\begin{equation}
    \ddot{c}^\mu + \Gamma^\mu_{\alpha\beta} \dot{c}^\alpha \dot{c}^\beta = 0
\end{equation}

onde $\Gamma^\mu_{\alpha\beta}$ s\~ao os s\'imbolos de Christoffel.

\subsection{Campo Direcional}

O campo direcional $\bm{F} \in \R^5$ captura a dire\c{c}\~ao preferencial de
evolu\c{c}\~ao no manifold, derivado da entropia dos padr\~oes de aten\c{c}\~ao:

\begin{equation}
    H_\ell = -\sum_t \bm{A}_{\ell,t} \log(\bm{A}_{\ell,t} + \epsilon)
\end{equation}

\begin{equation}
    \bm{F} = \text{MLP}_F(\bar{H}) \in \R^5
\end{equation}

onde $\bar{H}$ \'e a entropia m\'edia sobre camadas e heads.

\subsection{Curvatura de Ricci}

A sa\'ude do manifold pode ser avaliada pela curvatura escalar de Ricci,
computada a partir do tensor m\'etrico:

\begin{equation}
    R = g^{\mu\nu} R_{\mu\nu} = g^{\mu\nu} R^\lambda_{\ \mu\lambda\nu}
\end{equation}

Na implementa\c{c}\~ao, utilizamos uma aproxima\c{c}\~ao computacionalmente
eficiente baseada no \textit{condition number} de $\bG$:

\begin{equation}
    \text{health}_{\text{metric}} = 1 - \frac{\log(\kappa(\bG))}{\log(\kappa_{\max})}
\end{equation}

onde $\kappa(\bG) = \lambda_{\max} / \lambda_{\min}$ \'e o n\'umero de
condi\c{c}\~ao. Um manifold saud\'avel tem $\kappa(\bG) \approx 1$
(geometria isotr\'opica).

\subsection{Regulariza\c{c}\~ao do Tensor M\'etrico}

Para evitar degenera\c{c}\~ao da geometria, adicionamos uma regulariza\c{c}\~ao
no \textit{condition number}:

\begin{equation}
    \calL_{\text{metric}} = \log\big(\kappa(\bG) + 1\big)
\end{equation}

Isso for\c{c}a $\bG$ a permanecer pr\'oximo da identidade, evitando que
certas dire\c{c}\~oes do manifold colapsem ou se inflacionem
excessivamente. A forma logar\'itmica fornece uma penalidade mais suave
que $(\kappa - 1)^2$, permitindo pequenos desvios da isotropia enquanto
previne degenera\c{c}\~ao extrema.
